% ============================================================
%  BUFFER gate - logic diagram.
%  A buffer is simply TWO NOT gates in series:
%        buffer = NOT(NOT(A)) = A
%  Each NOT gate is the complementary (2N3904 + 2N3906) inverter
%  board from the `not` folder - nothing new to build at the
%  transistor level, just wire two of them together.
%  Compile:  pdflatex circuit.tex
% ============================================================
\documentclass[border=16pt]{standalone}
\usepackage{circuitikz}
\usepackage{amsmath}
\usetikzlibrary{calc}

\begin{document}
\begin{circuitikz}[line width=1.0pt, scale=1.6, transform shape]
  \ctikzset{logic ports=american, logic ports/thickness=1.1}

  % the two NOT gates
  \node[not port] (g1) at (0,0)   {};
  \node[not port] (g2) at (3.6,0) {};
  \node[below=3pt,font=\footnotesize] at (g1.south) {NOT};
  \node[below=3pt,font=\footnotesize] at (g2.south) {NOT};

  % input A
  \draw (g1.in) -- ++(-1.6,0) node[ocirc]{} node[left]{Input $A$};

  % middle node  A-bar
  \draw (g1.out) -- (g2.in);
  \node[above=1pt,font=\footnotesize] at ($(g1.out)!0.5!(g2.in)$) {$\overline{A}$};

  % output Y = A
  \draw (g2.out) -- ++(1.6,0) node[ocirc]{} node[right]{Output $Y=A$};

\end{circuitikz}
\end{document}
